\documentclass[11pt,a4paper]{article}
\usepackage[T1,T2A]{fontenc}
\usepackage[utf8]{inputenc}
\usepackage[russian]{babel}
\usepackage[margin=2.2cm]{geometry}
\usepackage{amsmath,amssymb,amsthm,mathtools}
\usepackage{enumitem}
\usepackage{booktabs,tabularx}
\usepackage{microtype}
\usepackage{xcolor}
\usepackage[colorlinks=true,linkcolor=blue!50!black,urlcolor=blue!50!black,citecolor=green!35!black]{hyperref}

\newtheorem{theorem}{Теорема}
\newtheorem{lemma}[theorem]{Лемма}
\newtheorem{proposition}[theorem]{Предложение}
\newtheorem{corollary}[theorem]{Следствие}
\theoremstyle{definition}
\newtheorem{remark}[theorem]{Замечание}
\newtheorem{example}[theorem]{Пример}
\newtheorem{question}{Вопрос}

\newcommand{\R}{\mathbb R}
\newcommand{\eps}{\varepsilon}
\DeclareMathOperator{\Hess}{Hess}
\DeclareMathOperator{\Sym}{Sym}
\DeclareMathOperator{\diag}{diag}
\setlist{itemsep=2pt,topsep=3pt}

\title{Проблема Арнольда 1971-4\\[3pt]
\large (обращение теоремы Лагранжа--Дирихле при изолированной критической точке):\\
что известно, какие идеи работают и где ломается двумерное доказательство}
\author{}
\date{3 октября 2026 г.}

\begin{document}
\maketitle
\vspace{-8mm}
\noindent\emph{Репозиторий доски: \nolinkurl{notes/vlad-claude-tg/}. Статусы утверждений по меткам доски
([Доказано], [Эскиз], [Литература], [Эксперимент]) сведены в \texttt{README.md} этой папки.}
\vspace{2mm}
\tableofcontents

\section*{Коротко}
\addcontentsline{toc}{section}{Коротко}
\begin{enumerate}[leftmargin=*]
\item Задача --- дословно проблема 1971-4 из книги <<Задачи Арнольда>>. При $n=2$ она решена:
Паламодов (1977, даже в усиленной <<тотальной>> форме), Талиаферро (1980); отдельное короткое
доказательство именно для изолированной критической точки дал Брюнелла (1998).
При $n\ge 3$ она \emph{открыта} --- насколько мне известно, и для полиномов от трёх переменных.
В работах 2022 и 2025 гг. прямо сказано, что даже эта <<облегчённая>> версия обращения
теоремы Лагранжа--Дирихле остаётся открытой; общее доказательство, намеченное Паламодовым в 1995 г.,
считается неполным, а <<короткое доказательство>> 2021 г. отозвано автором.
\item Почти все положительные результаты получаются одним из двух способов:
(а) построить \emph{асимптотическую траекторию}, входящую в равновесие;
(б) построить векторное поле $V$ с $\langle\nabla U,V\rangle\le\alpha U$ и $\Sym DV\ge-\frac\alpha2$;
тогда функция $\langle\dot q,V(q)\rangle$ монотонно растёт (вириальная функция Четаева--Паламодова).
В \S\ref{sec:virial} из (б) за несколько строк выводятся квазиоднородный случай, теорема Паламодова
о первой форме без критических точек и случай гессиана ранга $n-1$.
\item Где ломается $n=2$ (\S\ref{sec:2d}):
(i) у Брюнеллы ключевой объект --- \emph{трёхмерное} инвариантное множество $\{H=0\}$;
при $n=3$ оно пятимерно, а редукция особенностей векторных полей в размерности $\ge4$ неизвестна;
к тому же эйлерова характеристика линка равна $2\chi(\Sigma)$, $\Sigma=\{U\le0\}\cap S^2_\eps$,
и может обращаться в нуль;
(ii) в подходе через раздутия (Паламодов) на плоскости центры раздутий --- точки, в $\R^3$ --- кривые,
и теряется контроль над $\Sym DV$; ровно здесь пробел общего доказательства;
(iii) в подходе через асимптотические решения на плоскости множество <<кандидатов в асимптотические
направления>> конечно, а поперечная динамика одномерна; в $\R^3$ появляются \emph{жёлобы}~---
целые кривые вырожденных направлений --- и двумерная поперечная динамика.
\item Для полиномов от трёх переменных известные критерии сводят задачу к <<жёсткому ядру>>
(\S\ref{sec:3d}): $U=U_4+U_5+\dots$, где $U_4\ge0$ --- неотрицательная квартика с нулями на сфере
(по Гильберту и Безу: либо не более 4 нулей в $\mathbb{RP}^2$, либо $U_4=cQ^2$, либо $U_4=L^2P$),
плюс вырожденные подслучаи гессиана ранга~1. Там же --- тестовые примеры и возможные программы атаки.
\textbf{Обновление:} для примера~\ref{ex:hard}, не покрываемого известными критериями, тотальная неустойчивость
доказана в отдельной заметке <<Пример 13>> (\nolinkurl{notes/vlad-claude-tg/example13-proof.md}) явным вириальным полем, адаптированным к жёлобу. То же
рассуждение даёт общий критерий: жёлоб с \emph{регулярным} эффективным потенциалом.
\end{enumerate}

%=====================================================================
\section{Постановка и статус задачи}\label{sec:status}

\subsection{Постановка и два вида неустойчивости}
Рассматриваем $\ddot q=-\nabla U(q)$, $q\in\R^n$, с энергией $H=\tfrac12|\dot q|^2+U(q)$;
$U(0)=0$, $\nabla U(0)=0$, $0$ --- изолированная критическая точка, не минимум.
Тогда множество $\{U<0\}$ полуаналитично и примыкает к $0$; по лемме о выборе кривой в нём есть
аналитическая дуга, выходящая из $0$, поэтому $U$ принимает отрицательные значения на каждой
малой сфере $S_r=\{|q|=r\}$.

\begin{itemize}[leftmargin=*]
\item \emph{Асимптотическая траектория}: решение $q(t)\not\equiv0$ с $(q,\dot q)\to(0,0)$ при
$t\to+\infty$ или $t\to-\infty$. Её существование влечёт неустойчивость (лемма~\ref{lem:asympt}).
В.\,В.~Козлов (1982) высказал гипотезу, что у натуральной системы с аналитическим потенциалом
такая траектория существует всегда, когда $0$ не минимум.
\item \emph{Тотальная неустойчивость} (Паламодов): существует окрестность $\Omega$ точки $0$, такая что
\emph{любое} движение с $H<0$, стартовавшее в $\Omega$, покидает $\Omega$ за конечное время.
Если $0$ не минимум, отсюда следует неустойчивость по Ляпунову. Гипотеза Паламодова: всякая
неминимальная критическая точка аналитического потенциала тотально неустойчива.
\end{itemize}

\subsection{Что известно}
\paragraph{Гладкий класс.} Без аналитичности обращение неверно: пример Пенлеве--Уинтнера
$U(x)=e^{-1/x^2}\cos(1/x)$ при $n=1$, пример Лалуа (1976) при $n=2$; в обоих примерах критическая
точка \emph{не} изолирована. В $C^\infty$ устойчивость может зависеть от кинетической энергии
(Бертотти--Болотин, 2000). Пример Уреньи (2020): $C^\infty$-потенциал на плоскости со строгим
максимумом в $0$, $\nabla U\ne0$ при $q\ne0$, и с периодическими решениями сколь угодно близко к
равновесию. Методологически это важно: в гладком классе тотальная неустойчивость может нарушаться
даже при изолированной критической точке, так что любое доказательство тотальной
неустойчивости обязано существенно использовать аналитичность.

\paragraph{$n=2$: решено.} Паламодов [Pa77] доказал (даже тотальную) неустойчивость для любого
аналитического потенциала на плоскости, Талиаферро [Ta80b] независимо --- неустойчивость по Ляпунову;
Козлов [Ko81] --- для произвольной кинетической энергии; нестрогий минимум ---
Лалуа--Пайфер [LP82]. Для изолированной критической точки, не являющейся ни минимумом, ни
максимумом, Брюнелла [Br98] дал короткое концептуальное доказательство существования
асимптотических траекторий (\S\ref{sec:brunella}); случай максимума прост (Хагедорн [Ha71];
принцип Мопертюи--Якоби и теорема Хопфа--Ринова).

\paragraph{$n\ge3$: открыто.} См. введения к [BP22] и [Bu25] (<<Even this simplified version of the
Lagrange--Dirichlet converse remains open>>). Паламодов [Pa95] объявил теорему о тотальной
неустойчивости любой неминимальной критической точки, но доказательство неполно (это признаётся
в [BP22], [Bu21], [AR17]); заметка [Bu21] с <<коротким доказательством>> отозвана автором после того,
как Паламодов указал ошибку (\S\ref{sec:resolution}). Известные частичные результаты для $n\ge3$:
\begin{itemize}[leftmargin=*]
\item Ляпунов: гессиан имеет отрицательное собственное значение;
\item Четаев [Ch52]: $U$ однороден; Козлов--Паламодов [KP82] (также [Ko82], [MN89],
[Ta90]): если первая ненулевая однородная форма $U_k$ не имеет минимума, то существует асимптотическая
траектория;
\item Козлов [Ko86]: $U=u_2+u_m+\dots$, $u_2\ge0$, и $u_m$ не имеет минимума на $\ker u_2$ ---
асимптотическая траектория; Винер [Wi93]: коранг гессиана 2 (при технических условиях);
\item Паламодов [Pa20]: тотальная неустойчивость, если $U_k$ не имеет критических точек, кроме $0$
(теорема 3 там), если $U$ квазиоднороден, если старшая квазиоднородная часть не имеет критических
точек и хвост мал в специальном смысле (теорема 6), и редукция к ядру гессиана (теорема 5);
\item максимум: Хагедорн [Ha71], Талиаферро [Ta80a], Софер, Хагедорн--Мавен (вариационно);
\item Бургос [Bu25] (раздутие Мак-Ги): тотальная неустойчивость генерически
(если $U_k|_{S^{n-1}}$ --- функция Морса с регулярным значением $0$) и негенерический критерий
через вторую ненулевую однородную компоненту (\S\ref{sec:3d});
\item нестрогий минимум (не наш случай): [BMP21], [BP22].
\end{itemize}
Общая черта почти всех критериев: отсутствие минимума <<видно>> уже по первой ненулевой форме $U_k$
(иногда --- по квазиоднородной старшей части или по второй форме).
Ещё один полезный факт [KT18] (в изложении [Pa20]): для изолированного неустойчивого равновесия
аналитической системы \emph{общего вида} в чётномерном фазовом пространстве асимптотического
движения может не быть --- так что гипотеза Козлова существенно использует натуральную структуру.

\subsection{Про частный случай полиномов}\label{sec:poly}
Если $0$ --- изолированная критическая точка и в \emph{комплексной} области (число Милнора конечно),
росток $U$ аналитически конечно определён: в подходящих аналитических координатах $U$ совпадает со
своим многочленом Тейлора некоторой степени. Платой является замена евклидовой кинетической энергии
на аналитическую риманову метрику $g(q)$. Большинство известных методов (Козлов, Брюнелла,
вириальные поля) к виду аналитической метрики нечувствительны, поэтому в этом случае
<<полиномиальность>> сути не упрощает. Если же точка изолирована лишь над $\R$ (как в примерах
с $U_4=(x^2+y^2-z^2)^2$ ниже), то по теореме Тужрона росток лишь $C^\infty$-конечно определён,
и редукция идёт через гладкие (не аналитические) замены. Зато полиномиальность даёт
\emph{вычислимость}: многогранники Ньютона, полуалгебраическую геометрию, классификацию
неотрицательных форм (\S\ref{sec:3d}).

%=====================================================================
\section{Вириальный метод и что им доказывается}\label{sec:virial}

\begin{lemma}\label{lem:asympt}
Если существует решение $q(t)\not\equiv0$ с $(q(t),\dot q(t))\to(0,0)$ при $t\to-\infty$ или при
$t\to+\infty$, то равновесие неустойчиво.
\end{lemma}
\begin{proof}
Система обратима ($q(-t)$ --- тоже решение), поэтому достаточно случая $t\to-\infty$.
Положим $\eps=\frac12|(q(0),\dot q(0))|>0$. Если бы равновесие было устойчиво, нашлось бы $\delta>0$,
такое что траектории из $\delta$-окрестности не выходят из $\eps$-окрестности; но
$(q(t_0),\dot q(t_0))$ лежит в $\delta$-окрестности при $t_0\ll0$, а в момент $0$ траектория вне
$\eps$-окрестности.
\end{proof}

\begin{lemma}[вириальная лемма; по существу теорема 2 из {[Pa77]}, см. {[Pa20]}]\label{lem:virial}
Пусть $B$ --- шар с центром $0$, $V$ --- ограниченное $C^1$-поле на $B\cap\{U<0\}$ и при некотором
$\alpha>0$
\[
\text{(i)}\ \ \langle\nabla U,V\rangle\le\alpha\,U,\qquad
\text{(ii)}\ \ \langle DV(q)\,\xi,\xi\rangle\ge-\tfrac{\alpha}{2}|\xi|^2\quad\forall\,\xi\in\R^n .
\]
Тогда любая траектория с энергией $h<0$, стартовавшая в $B$, покидает $B$ не позже чем через время
$2M/(\alpha|h|)$, где $M=\sup|V|\cdot(2\sup_B|U|)^{1/2}$. В частности, если $0$ не минимум,
равновесие тотально неустойчиво.
\end{lemma}
\begin{proof}
На траектории $U(q)=h-\frac12|\dot q|^2\le h<0$, так что $V$ определено, пока она в $B$.
Для $F(t)=\langle\dot q,V(q)\rangle$ имеем $|F|\le M$ и
\[
\dot F=-\langle\nabla U,V\rangle+\langle DV\,\dot q,\dot q\rangle\ \ge\ -\alpha U-\tfrac\alpha2|\dot q|^2
=-\alpha h=\alpha|h|>0 .
\]
Значит, траектория проводит в $B$ время не больше $2M/(\alpha|h|)$.
\end{proof}

\begin{remark}
Для лагранжиана $\frac12 g_q(\dot q,\dot q)-U$ верно тождество
$\frac{d}{dt}\,g(\dot q,V)=\frac12(L_Vg)(\dot q,\dot q)-dU(V)$, и условие (ii) заменяется на
$L_Vg\ge-\alpha g$. У Паламодова вместо (ii) требуется неотрицательность формы $\langle DV(0)\xi,\xi\rangle$
в самой точке $0$ плюс $C^1$-гладкость поля вплоть до $0$.
\end{remark}

\begin{corollary}[квазиоднородный потенциал; {[Pa20, \S5]}]\label{cor:qh}
Если $U(\lambda^{a_1}q_1,\dots,\lambda^{a_n}q_n)=\lambda^dU(q)$, $a_i,d>0$, и $U$ принимает отрицательные
значения, то равновесие тотально неустойчиво.
\end{corollary}
\begin{proof}
$V=\frac1d\sum a_iq_i\partial_i$: по Эйлеру $\langle\nabla U,V\rangle=U$, $DV=\diag(a_i/d)\ge0$.
\end{proof}

\begin{example}
$U=x^4+y^4-z^6$ (веса $3,3,2$, $d=12$). В [BP22] он приведён как пример, к которому не применим ни один
критерий <<по первой ненулевой форме>> ($x^4+y^4\ge0$), но функция $3x\dot x+3y\dot y+2z\dot z$ сразу даёт
тотальную неустойчивость.
\end{example}

\begin{corollary}[Паламодов {[Pa20, теорема 3]}]\label{cor:pal}
Пусть $U=U_k+U_{k+1}+\dots$ --- разложение на однородные формы, $\nabla U_k\ne0$ на $S^{n-1}$ (т.е.
$0$ --- регулярное значение $U_k|_{S^{n-1}}$), и $0$ --- не минимум. Тогда равновесие тотально
неустойчиво.
\end{corollary}
\begin{proof}[Доказательство (в духе {[Pa20]})]
Из однородности $|\nabla U(q)|\ge c|q|^{k-1}$ вблизи $0$. Положим
\[
g=U-\tfrac1k\langle q,\nabla U\rangle=\sum_{j>k}\Bigl(1-\frac jk\Bigr)U_j=O(|q|^{k+1}),\qquad
V=\frac qk+g\,\frac{\nabla U}{|\nabla U|^2}.
\]
Тогда $\langle\nabla U,V\rangle=U$ точно, $|V|\le|q|/k+C|q|^2$, а оценки $\nabla g=O(|q|^k)$,
$\Hess U=O(|q|^{k-2})$ дают $D\bigl(g\nabla U/|\nabla U|^2\bigr)=O(|q|)$. Итак, $DV=\frac1kI+O(|q|)\ge0$
вблизи $0$, и работает лемма~\ref{lem:virial} с $\alpha=1$. (При этом условии из $U_k\ge0$ следовало бы
$U_k>0$ на сфере, т.е. строгий минимум; так что <<не минимум>> здесь означает, что $U_k$ принимает
отрицательные значения.)
\end{proof}

\begin{corollary}[гессиан ранга $n-1$; частный случай теоремы 5 из {[Pa20]}]\label{cor:rank}
Пусть $\Hess U(0)\ge0$ имеет одномерное ядро и $0$ --- не минимум. Тогда равновесие тотально неустойчиво.
\end{corollary}
\begin{proof}
Координаты $q=(x,y)$, $x\in\R$ вдоль ядра, $y\in\R^{n-1}$, $\partial^2_{yy}U(0)>0$. По теореме о неявной
функции $\partial_yU(x,\varphi(x))=0$, $\varphi(0)=\varphi'(0)=0$. Приведённый потенциал $u(x)=U(x,\varphi(x))$;
при $w=y-\varphi(x)$ имеем $U=u(x)+\frac12\langle A(x)w,w\rangle+O(|w|^3)$, $A>0$. Поэтому $0$ не минимум
для $U$ $\iff$ не минимум для $u$, и $u=cx^m(1+O(x))$, $m\ge3$. Положим
\[
V_x=\frac{u(x)}{u'(x)}=\frac xm+O(x^2),\qquad V_y=\varphi'(x)V_x+\beta\,w,\qquad \beta=\tfrac14 .
\]
Дифференцируя тождество $\partial_yU(x,\varphi(x))=0$, получаем $\partial_x\partial_yU(x,\varphi)=-A\varphi'$;
отсюда перекрёстные члены сокращаются и
\[
\langle\nabla U,V\rangle=u+\beta\langle Aw,w\rangle+O\bigl(|w|^2(|x|+|w|)\bigr)\ \le\ U
\]
вблизи нуля (так как $\beta<\frac12$). Далее $DV=\begin{pmatrix}1/m+O(x)&0\\ O(x)&\beta I\end{pmatrix}$, т.е.
$\Sym DV>0$, и работает лемма~\ref{lem:virial} с $\alpha=1$.
\end{proof}

\begin{remark}[квазиоднородная старшая часть и анизотропия]
Естественное обобщение следствия~\ref{cor:pal} на случай, когда старшая часть $u$ квазиоднородна
с весами $w_i$ и не имеет критических точек вне нуля, требует поправки к взвешенному полю Эйлера.
Её производная $\partial_jV_i$ имеет взвешенный порядок $w_i-w_j+\delta$, что при $w_j>w_i+\delta$
\emph{не мало} --- проявление того же эффекта, что и в \S\ref{sec:resolution}. Поэтому в теореме~6
из [Pa20] на хвост $v=U-u$ наложено более сильное условие $v=O(|x|\,|x|_w)$.
\end{remark}

\subsection{Почему вириальной леммы не хватает}\label{sec:limits}
\begin{enumerate}[leftmargin=*]
\item \emph{Теорема Сайто} (1971): для ростка с изолированной особенностью аналитическое в нуле поле $V$
с $V(U)=U$ существует тогда и только тогда, когда $U$ квазиоднороден в подходящих координатах.
Поэтому в общем случае поле приходится строить <<руками>>, и оно не аналитично в нуле
(в следствии~\ref{cor:pal} оно лишь $C^1$).
\item \emph{Каноническое поле} $V_*=U\nabla U/|\nabla U|^2$ определено вне нуля благодаря
изолированности, $V_*(U)=U$, и $|V_*|\le C|q|$ по неравенству Бохнака--Лоясевича $|q||\nabla U|\ge c|U|$.
Но на векторах $\xi$, касательных к поверхности уровня, $\langle DV_*\xi,\xi\rangle=U\langle\Hess U\,\xi,\xi\rangle/|\nabla U|^2$
--- это определяется кривизной поверхностей уровня и может быть сколь угодно отрицательным.
Уже для $U=-x^2+3y^2$ в точках $(x,0)$ получаем $\partial_yV_*^y=-\frac32<-\frac12$, хотя поле Эйлера
$q/2$ годится. Можно добавлять любые поля $W$, касательные к уровням ($W(U)=0$); вопрос в том, как выбрать
$W$, чтобы $\Sym D(V_*+W)\ge-\frac12$, --- это глобальная задача, и общего рецепта нет.
\item Если гипотеза Паламодова о тотальной неустойчивости неверна для какого-то аналитического $U$, то
поля из леммы~\ref{lem:virial} для него не существует вовсе, и доказывать неустойчивость пришлось бы
через асимптотические траектории (лемма~\ref{lem:asympt}).
\end{enumerate}

%=====================================================================
\section{Как доказывается случай $n=2$ и где ломается $n=3$}\label{sec:2d}

\subsection{Метод Брюнеллы}\label{sec:brunella}
\begin{theorem}[Брюнелла, 1998]
Пусть $v$ --- аналитическое векторное поле в окрестности $0\in\R^N$ с изолированной особой точкой $0$,
$M$ --- неприводимое трёхмерное аналитическое подмножество, содержащее $0$, гладкое вне $0$ и
инвариантное относительно $v$, причём линк $M$ в нуле имеет ненулевую эйлерову характеристику. Тогда
найдётся траектория в $M\setminus\{0\}$, стремящаяся к $0$ при $t\to+\infty$ или при $t\to-\infty$.
\end{theorem}

\paragraph{Применение при $n=2$.} $N=4$, $v=X_H$, $M$ --- компонента $\{H=0\}$. Гладкость $M$ вне нуля ---
\emph{ровно то место, где используется изолированность} критической точки: $dH=(\nabla U,\dot q)\ne0$
на $M\setminus\{0\}$. Над каждой дугой множества $\Sigma=\{U\le0\}\cap S^1_\eps$ слой
$\{|p|=\sqrt{-2U(q)}\}$ --- окружность, стягивающаяся в точку на концах дуги; поэтому каждая компонента
линка --- сфера $S^2$, $\chi=2\ne0$. (Для максимума $\Sigma=S^1$, линк --- тор, $\chi=0$; но этот случай
прост.) Далее лемма~\ref{lem:asympt}.

\paragraph{Ингредиенты доказательства:} (a) локальная десингуляризация трёхмерных векторных полей
(Кано) и эквидесингуляризация (Санчо де Салас); (b) формула Пуанкаре--Хопфа для поля, индуцированного на
исключительном дивизоре раздутия ($\cong$ линку); (c) разбор элементарных особых точек в размерности~3:
центральные многообразия размерности $\le2$ и планарная теория.

\paragraph{Что ломается при $n=3$.}
\begin{enumerate}[leftmargin=*]
\item $\dim M=2n-1=5$. Редукция особенностей векторных полей известна в размерности 2 (Бендиксон,
Зайденберг) и 3 (Кано; Панаццоло 2006; Маккуиллан--Панаццоло), а в размерности $\ge4$, насколько мне
известно, это открытая проблема.
\item Даже после гипотетической редукции элементарные особые точки в размерности~5 имеют центральные
многообразия размерности до~4, на которых динамика снова вырождена и никакой классификации нет.
\item Индексное условие может нарушаться:
\end{enumerate}

\begin{proposition}\label{prop:chi}
Пусть $\Sigma=\{U\le0\}\cap S^{n-1}_\eps$ ($\eps$ мало, $0$ --- регулярное значение $U|_{S_\eps}$) и
$L$ --- линк множества $\{H=0\}$ в нуле. Тогда
\[
\chi(L)=\chi(S^{n-1})\bigl(\chi(\Sigma)-\chi(\partial\Sigma)\bigr)+\chi(\partial\Sigma),
\qquad\text{в частности}\quad \chi(L)=2\chi(\Sigma)\ \ \text{при } n=3 .
\]
\end{proposition}
\begin{proof}[Набросок]
На $\{H=0\}$ функция $|q|^2$ --- <<ковровая>>, так что $L\cong\{|q|=\eps,\ U(q)\le0,\ |p|^2=-2U(q)\}$.
Проекция на $q$ имеет слой $S^{n-1}$ над внутренностью $\Sigma$ и точку над $\partial\Sigma$; остаётся
аддитивность эйлеровой характеристики с компактными носителями. При $n=3$ граница $\partial\Sigma$ ---
объединение окружностей, $\chi(\partial\Sigma)=0$.
\end{proof}
При $n=3$ область $\Sigma\subset S^2$ планарна, и $\chi(\Sigma)=0$ ровно для колец. Если все компоненты
$\Sigma$ --- кольца (например, $U=z^2-(x^2+y^2)^2$: отрицательная область --- тонкий поясок вокруг экватора),
индексное рассуждение ничего не даёт --- аналог <<плохого>> случая максимума при $n=2$, но уже не
тривиальный. (Конкретно этот пример квазиоднороден и покрывается следствием~\ref{cor:qh}.)
Итак, двумерная теорема Брюнеллы --- это на самом деле теорема о векторных полях в размерности~3;
для $n=3$ понадобился бы её аналог для пятимерных многообразий, что выглядит очень трудным.

\subsection{Вириальное поле через раздутия (Паламодов) и пробел при $n\ge3$}\label{sec:resolution}
Для $n=2$ тотальная неустойчивость установлена в [Pa77]; согласно [Pa20], там же впервые появился
критерий леммы~\ref{lem:virial}. В [Pa95] предложена схема построения такого поля для любого $n$
(индукция по раздутиям с поправками поля на каждом шаге); её упрощённый вариант из [Bu21] таков: раздуть $0$, затем принципализовать идеал
$(U)$ по Хиронаке, $\sigma\colon\widetilde M\to\R^n$; в окрестности каждой точки $p$ дивизора над нулём
есть координаты $w$ с $\widetilde U=\pm w_1^{d_1}\cdots w_n^{d_n}$; поле Эйлера
$\widetilde V_p=\sum w_i\partial_{w_i}$ удовлетворяет $\widetilde V_p(\widetilde U)=c_p\widetilde U$; затем
$V_p=\sigma_*\widetilde V_p$ и $V=\sum f_pV_p$ (разбиение единицы). Условие (i) выполняется
автоматически: $V(U)=PU$, $P=\sum f_pc_p\ge1$. Вся трудность --- в условии (ii).
\begin{enumerate}[leftmargin=*]
\item \emph{Образ поля Эйлера при раздутии не <<растягивающий>>.} Пример: раздуем $0\in\R^2$ (карта
$x=a$, $y=aw$), затем точку $(a,w)=(0,c)$: $x=a$, $y=ca+a^2b$. Тогда
\[
\sigma_*(a\partial_a+b\partial_b)=x\partial_x+(3y-2cx)\partial_y,\qquad
\Sym DV=\begin{pmatrix}1&-c\\-c&3\end{pmatrix},
\]
что знаконеопределено при $|c|>\sqrt3$; в цепочке из $j$ таких раздутий внедиагональный член растёт
как $jc$. Здесь $c$ убирается поворотом, но при раздутиях вдоль \emph{кривых} (а при $n=3$ без них не
обойтись: см. жёлобы в \S\ref{sec:puiseux}) аналогичный параметр меняется вдоль центра, и глобально
<<повернуть>> его нельзя.
\item \emph{Склейка.} В $\langle DV\xi,\xi\rangle$ входит
$\sum_p\xi(f_p)\langle\xi,V_p\rangle=\sum_p\xi(f_p)\langle\xi,V_p-V_{p'}\rangle$; срезки меняются на
масштабе тонких <<рогов>> $\sigma(W_p)$, а разность $V_p-V_{p'}$ на пересечении, вообще говоря, того же
порядка, что ширина рога, --- получается член порядка $|\xi|^2$ без малости (это эвристика, а не
теорема, но именно такой член оказался роковым, см. ниже).
\end{enumerate}
В отозванной заметке [Bu21] утверждалось $\langle\xi,\nabla_\xi V\rangle=(1+o(1))|\xi|^2$. По словам самого
автора, <<второй член последнего равенства леммы 2.2 не имеет нужной асимптотики>>; в последней версии
заметки это равенство для склеенного поля, и его второй член --- склеивающее слагаемое
$\sum_p\xi(f_p)\langle\xi,V_p\rangle$. Кроме того, уже для отдельного $V_p$ там отброшено слагаемое
$\frac12V_p(g_{ab})\xi^a\xi^b$ с производной метрики, записанной в сингулярных координатах раздутия; оно
того же порядка, что $|\xi|^2$ (ср. пример выше). Итог: нужна <<разрешающая>> процедура,
\emph{совместимая с механикой}, т.е. с количественным контролем того, как спуск через раздутия искажает
симметричную часть производной поля. На плоскости центры раздутий --- точки и дивизоры --- кривые;
в $\R^3$ центры --- кривые, дивизоры --- поверхности, и такого контроля пока нет.

\subsection{Асимптотические решения и разложения Пюизё}\label{sec:puiseux}
Двумерные рассуждения через асимптотические траектории (Козлов; Талиаферро; Гарсия--Таль) так или иначе
опираются на специфически двумерные факты.
\begin{enumerate}[leftmargin=*,label=(F\arabic*)]
\item Нулевое множество $\{U=0\}$ и полярная кривая $\{\partial_\varphi U=0\}$ вблизи нуля --- конечные
объединения полуветвей Пюизё; $\{U<0\}$ --- конечное объединение криволинейных секторов.
\item \emph{Кандидатов в асимптотические направления конечное число.} Асимптотическая траектория подходит
к множеству критических точек $f=U_k|_{S^{n-1}}$ с $f\le0$ ([Bu25, предл.~1.5]; на языке раздутия
Мак-Ги --- к равновесиям на <<коллизионном>> многообразии). При $n=2$ и $U_k\ge0$ это нули бинарной формы
--- конечное множество. Схематически: каждая вырожденная точка разрешается конечным числом взвешенных раздутий
(алгоритм Ньютона--Пюизё), в конце --- квазиоднородное укорочение с <<лучевым>> решением, а затем
настоящее решение (метод Козлова; теорема Кузнецова: формальное решение $\Rightarrow$ настоящее).
\item \emph{Поперечная структура одномерна}: ограничение $U$ на дугу окружности --- функция одной
переменной, её минимум имеет конечный порядок, и алгоритм Ньютона--Пюизё обрывается.
\item Топология плоскости (теорема Жордана, промежуточные значения) для выбора нужной траектории.
\end{enumerate}
При $n=3$:
\begin{enumerate}[leftmargin=*,label=(F\arabic*$'$)]
\item $\Sigma$ --- планарная область на $S^2$ любой топологии; $\{U=0\}$ --- поверхность, возможно
с неизолированными особенностями (конус над кривой); вместо Пюизё --- разрешение особенностей
поверхностей (Зарисский, Абхьянкар, Хиронака) с центрами-кривыми.
\item \emph{Жёлобы.} Множество кандидатов может быть \emph{кривой}: например, $U_4=(x^2+y^2-z^2)^2$
обращается в нуль на двух окружностях сферы. На границе раздутия Мак-Ги получается целая кривая
негиперболических равновесий; выбор направления вдоль жёлоба определяется <<эффективным потенциалом>>
--- следующими членами разложения, ограниченными на жёлоб, --- который сам может быть вырожден (каскад
раздутий вдоль кривых без априорной конечности). Возникает и вопрос, отсутствующий при $n=2$: сходится ли
асимптотическая траектория к \emph{точке} жёлоба (аналог гипотезы Тома о градиенте, доказанной для
градиентных потоков Курдыкой--Мостовским--Парусинским)?
\item \emph{Поперечная структура двумерна}: минимум функции двух переменных может быть вырожден по
Ньютону в любых координатах; поперечная динамика --- осциллятор с двумя степенями свободы и медленно
меняющимися параметрами: резонансы, поворот главных осей (теорема Реллиха), обмен энергией; одно
квазиоднородное укорочение не схватывает все масштабы.
\item Аргументы типа теоремы Жордана недоступны; нужны принцип Важевского или индекс Конли с
нетривиальной топологией множества выхода.
\end{enumerate}

\subsection{Вариационный подход и почему он не работает <<в лоб>>}
Для максимума (и вообще при $U\le0$) работает принцип Мопертюи--Якоби: при $h>0$ метрика
$(h-U)|dq|^2$ полна и невырождена, по Хопфу--Ринову есть геодезическая из $0$ в любую точку, отсюда
неустойчивость. Для седловых $U$ граница области Хилла $\{U=h\}$ --- место, где метрика Якоби
\emph{наименьшая}: минимизирующие кривые <<прилипают>> к ней (тормозные орбиты Болотина), и подобласть
$\{U\le c\}$ оказывается вогнутой в метрике Якоби (геодезические изгибаются в сторону меньших $U$,
т.е. внутрь области, и могут касаться её границы изнутри). Нужна область, выпуклая в метрике Якоби и
<<выводящая>> из окрестности нуля, --- её построение по трудности эквивалентно исходной задаче.

%=====================================================================
\section{Полиномы от трёх переменных}\label{sec:3d}

\subsection{Сведение к <<жёсткому ядру>>}
Пусть $U$ --- полином (или аналитическая функция) в $\R^3$ с изолированной неминимальной критической
точкой в нуле.
\begin{enumerate}[leftmargin=*]
\item Гессиан имеет отрицательное собственное значение --- Ляпунов. Иначе $\Hess U(0)\ge0$.
\item Ранг $2$ --- следствие~\ref{cor:rank} (тотальная неустойчивость).
\item Ранг $1$ (коранг 2): задача для приведённого двумерного потенциала $u$ на ядре. Если первая форма
на ядре не имеет минимума --- Козлов [Ko86]; если старшая форма $u$ без критических точек или $u$
квазиоднороден --- [Pa20, теорема 5] и замечание после следствия~\ref{cor:rank}; много случаев покрывает
Винер [Wi93]. В общем случае естественно ожидать сведения к теореме для $n=2$, но строго это не сделано
(медленное многообразие $y=\varphi(x)$ инвариантно лишь приближённо).
\item $\Hess U(0)=0$. Если $U_3\not\equiv0$, кубическая форма принимает отрицательные значения ---
Козлов--Паламодов [KP82]; если $U_3\equiv0$ и $U_4$ принимает отрицательные значения --- то же.
\item Остаётся \textbf{жёсткое ядро}: $U=U_4+U_5+\dots$, $U_4\ge0$ и $U_4$ обращается в нуль на сфере
(при $U_4>0$ на сфере был бы строгий минимум); аналогично при старшей форме чётной степени $\ge6$.
Если вдобавок $\nabla U_4\ne0$ вне нуля, то $U_4>0$ --- противоречие, так что $0$ --- критическое
значение $U_4|_{S^2}$, и общие генерические критерии [Bu25] не применяются.
\end{enumerate}

\subsection{Неотрицательные тернарные квартики и жёлобы}
По теореме Гильберта (1888) неотрицательная тернарная квартика есть сумма трёх квадратов квадратичных
форм $q_1^2+q_2^2+q_3^2$; её нули --- общие нули коник $q_i$. По Безу отсюда получается классификация:
\begin{itemize}[leftmargin=*]
\item либо нулей в $\mathbb{RP}^2$ конечное число, и их не больше четырёх
(например, $(x^2-z^2)^2+(y^2-z^2)^2$ с нулями $(\pm1,\pm1,1)$);
\item либо нули образуют кривую, и тогда $U_4=cQ^2$ ($Q$ --- знаконеопределённая квадратичная форма;
жёлоб --- коника, в т.ч. пара прямых при $Q=L_1L_2$) или $U_4=L^2P$ ($P\ge0$; жёлоб --- прямая $L=0$).
\end{itemize}
В первом случае около каждого нулевого направления возникает задача с двумерной поперечной структурой
(часто сводимая к квазиоднородным укорочениям), во втором --- жёлоб, принципиально новое по сравнению
с $n=2$ явление. Критерий Бургоса [Bu25, следствие~1.4] в ньютоновском случае: если вторая ненулевая
однородная компонента $U_{l_2}$ \emph{отрицательна во всех} точках множества
$\{\text{крит. точки } U_{l_1}|_{S^2}\ \text{с}\ U_{l_1}\le0\}$ (для жёсткого ядра --- во всех точках жёлоба),
то равновесие тотально неустойчиво. Для $U=(x^2+y^2-z^2)^2+W$, $W=W_m+\dots$, $m\ge5$, это условие можно
ослабить (заметка <<Пример 13>>, предложение 1): достаточно, чтобы $0$ было регулярным значением ограничения
$W_m$ на жёлоб и чтобы это ограничение принимало отрицательные значения.

\subsection{Тестовые примеры}
\begin{example}[покрыт]
$U=(x^2+y^2-z^2)^2-z^6$. Критическая точка изолирована (проверка: при $Q=x^2+y^2-z^2\ne0$ из
$xQ=yQ=0$ следует $x=y=0$, далее $z=0$; при $Q=0$ из третьего уравнения $z=0$). Здесь $U_6=-z^6<0$
на жёлобе $\{Q=0\}\cap S^2$ --- работает критерий Бургоса. Кроме того, $U$ инвариантен относительно
вращений вокруг оси $z$, плоскость $\{y=0\}$ инвариантна для динамики, и задача сводится к $n=2$.
\end{example}

\begin{example}[покрыт через симметрию]
$U=(x^2+y^2-z^2)^2+xz^5$. Критическая точка изолирована (аналогичная выкладка), $U_6=xz^5$ меняет знак на
жёлобе, так что критерий Бургоса не работает. Но $U$ чётен по $y$: плоскость $\{y=0\}$ инвариантна
(при $y=\dot y=0$), а $U(x,0,z)=(x^2-z^2)^2+xz^5$ не имеет минимума (на прямой $x=-z$ он равен $-z^6$).
По двумерной теореме равновесие неустойчиво. Мораль: любая изометрическая симметрия с неподвижным
подпространством $W$, на котором $U|_W$ не имеет минимума, сводит размерность.
\end{example}

\begin{example}[не покрыт известными критериями; доказан в заметке <<Пример 13>>]\label{ex:hard}
\[
U=(x^2+y^2-z^2)^2+z^5x+z^4xy+\tfrac12 z^6 .
\]
\begin{itemize}[leftmargin=*]
\item \emph{Изолированность.} При $z=0$ из $4xQ=4yQ=0$ сразу $x=y=0$. При $z\ne0$, $x=az$, $y=bz$, два
уравнения критичности становятся однородными и дают $a=-3b/((4b+1)(b+1))$ и
$(1+b)^3(4b+1)^2=9b$; единственный вещественный корень $b\approx-1.7577$, $a\approx1.1540$, и
$a^2+b^2\approx4.42\ne1$, тогда как третье уравнение требует $a^2+b^2-1=O(z^2)$. Значит, вблизи нуля
других критических точек нет (численная проверка на сферах согласуется).
\item \emph{Не минимум, не максимум}; $\Hess U(0)=0$, $U_3=0$, $U_4=Q^2\ge0$ с жёлобом из двух окружностей
$\{z=\pm r/\sqrt2\}$; $\nabla U_4=0$ на жёлобе.
\item $U_6=z^5x+z^4xy+\frac12z^6$ на верхней окружности жёлоба равен
$\frac{r^6}{8}G(\varphi)$, $G(\varphi)=\cos\varphi+\frac12\sin2\varphi+\frac12$, и меняет знак, так что
следствие~1.4 из [Bu25] не применимо. Ни Ляпунов, ни Козлов--Паламодов, ни [Pa20] не применимы; $U$ не
квазиоднороден в данных координатах; его единственная изометрическая симметрия --- $q\mapsto-q$, без
инвариантных плоскостей и прямых.
\end{itemize}
\emph{Эвристика.} Вблизи верхней полости конуса при $\eta=\sqrt{x^2+y^2}-z$ имеем
$Q\approx2z\eta$ и $U\approx4z^2\eta^2+\frac{r^6}{8}G(\varphi)$: поперечная жёсткость $\sim r^2$,
эффективный потенциал вдоль жёлоба $\sim r^6G(\varphi)$. У $G$ невырожденный отрицательный минимум
$G(5\pi/6)\approx-0.799$. Ожидаемая асимптотическая траектория: $\varphi\to5\pi/6$, $r\sim|t|^{-1/2}$,
поперечное смещение $\eta\sim r^3$ (баланс $8z^2\eta\sim\partial_\eta U_6\sim r^5$), что много меньше
поперечного масштаба $r^2$.

\emph{Обновление: доказано.} В координатах $s=(\rho+z)/\sqrt2$, $d=(\rho-z)/\sqrt2$ (где $\rho=\sqrt{x^2+y^2}$)
точно $Q^2=4s^2d^2$, и поле
$V=\frac16(s\partial_s+2d\partial_d)-\kappa s^{-2}(\partial_dU)\partial_d-\kappa' s^{-6}(\partial_\varphi U)\partial_\varphi$
($\kappa=\frac1{100}$, $\kappa'=\frac1{50}$) удовлетворяет условиям леммы~\ref{lem:virial} с $\alpha=\frac12$ на
$\{U<0\}$ вблизи нуля. Отсюда тотальная неустойчивость. Ключевое свойство: у $G$ нет кратных нулей.
Подробности --- в отдельной заметке.
\end{example}

\begin{example}[жёлоб, который сам надо искать]
$U=(Q-z\,m(q))^2-z^8$ с квадратичной формой $m$. Здесь $U_4=Q^2$, а $U_5=-2Qzm$ и $U_6=z^2m^2$ на жёлобе
дают нулевой и неотрицательный эффективные потенциалы (так что и критерий Бургоса не работает:
$U_{l_2}=U_5=0$ на жёлобе), а отрицательность живёт на \emph{деформированном} конусе $\{Q=zm\}$,
где $U=-z^8$. Это трёхмерный аналог двумерного примера $(y-x^2)^2-x^6$: правильный
жёлоб находится только процедурой типа Ньютона--Пюизё (выпрямление поверхности, а не кривой).
Изолированность критической точки надо проверять отдельно при конкретном выборе $m$.
\end{example}

\subsection{Возможные программы и конкретные вопросы}
\begin{enumerate}[leftmargin=*,label=(P\arabic*)]
\item \emph{Раздутия + нормальная гиперболичность.} Раздутие Мак-Ги (или квазиоднородное) равновесия;
в жёстком ядре на границе возникают кривые негиперболических равновесий (жёлобы), затем раздутия
вдоль них в духе Дюмортье--Руссари и Крупы--Шмоляна, инвариантные многообразия Фенихеля и медленные
потоки, управляемые эффективными потенциалами. Геометрическая часть доступна (разрешение особенностей
поверхностей в $\R^3$ классическое); нужна мера сложности, убывающая при раздутиях и совместимая
с динамикой.
\item \emph{Формальные решения + теорема Кузнецова.} Свести задачу к существованию \emph{формального}
асимптотического решения (ряды по $t^{-1/N}$, возможно с логарифмами) --- чисто алгебро-комбинаторная
задача степенной геометрии Брюно для системы ОДУ. При $n=2$ его даёт многоугольник Ньютона; при $n=3$
возможное препятствие --- движение вдоль жёлоба, не описываемое степенными рядами.
\item \emph{Индексная теория в размерности 5 для натуральных систем.} Использовать дополнительную
структуру (обратимость $(q,p)\mapsto(q,-p)$, квадратичность по импульсам), чтобы обойти отсутствие общей
редукции особенностей; индекс $\chi(L)=2\chi(\Sigma)$ подсказывает, что случаи с кольцевыми $\Sigma$
придётся разбирать отдельно.
\item \emph{Вириальные поля через торические раздутия.} Для потенциалов, невырожденных по Ньютону
(Кушниренко), разрешение торическое и явное; если $U$ невырожден и $0$ не минимум, то (с естественными
оговорками про координатные подпространства) некоторое квазиоднородное укорочение по грани многогранника
Ньютона принимает отрицательные значения --- по лемме о выборе кривой и невырожденности. Вопрос ---
склеить взвешенные поля Эйлера разных граней с контролем $\Sym DV$ (ср. условие $v=O(|x||x|_w)$ в [Pa20]).
\item \emph{Техника Курдыки--Мостовского--Парусинского} (характеристические показатели, контроль $U/|q|^\ell$
вдоль траекторий) как источник функций Четаева в <<рогах>> вокруг жёлобов.
\end{enumerate}

\begin{question} Верна ли неустойчивость для всех потенциалов в $\R^3$, невырожденных по Ньютону, с изолированной
неминимальной критической точкой?\end{question}
\begin{question} Пусть $U=Q^2+U_6+\dots$, $Q$ --- знаконеопределённая квадратичная форма. Если ограничение $U_6$
на жёлоб имеет только простые нули, то для $Q=x^2+y^2-z^2$ ответ положителен (заметка <<Пример 13>>).
Что происходит при кратных нулях эффективного потенциала, для неосесимметричного конуса, для особого жёлоба
или когда жёлоб надо сначала <<искать>> (пример~14)?\end{question}
\begin{question} Имеет ли каждая асимптотическая траектория натуральной системы с аналитическим потенциалом
предельное направление (аналог гипотезы Тома)?\end{question}
\begin{question} Всегда ли у $\ddot q=-\nabla U$ при неминимальной изолированной критической точке есть формальное
степенно-логарифмическое асимптотическое решение?\end{question}

%=====================================================================
\section{Какие разделы математики здесь работают}
\begin{center}
\small
\begin{tabularx}{\textwidth}{>{\raggedright\arraybackslash}p{4.3cm}X}
\toprule
Раздел & Инструменты и их роль\\
\midrule
Теория устойчивости & функции Ляпунова--Четаева, вириальные тождества, лемма~\ref{lem:virial}; тотальная неустойчивость\\
Асимптотические методы & квазиоднородные укорочения, степенная геометрия Брюно, метод Козлова--Фурты, теорема Кузнецова (формальное $\Rightarrow$ настоящее)\\
Вещественная аналитическая и полуалгебраическая геометрия & лемма о выборе кривой, неравенства Лоясевича и Бохнака--Лоясевича, конечная определённость (Мазер, Тужрон), разрешение особенностей (Хиронака; для поверхностей Зарисский--Абхьянкар), многогранники Ньютона и невырожденность по Кушниренко, о-минимальность, гипотеза Тома о градиенте (KMP)\\
Коммутативная алгебра & теорема Сайто ($U\in J(U)\iff$ квазиоднородность), целозамкнутость якобиева идеала\\
Особенности векторных полей & раздутия, редукция особенностей (Зайденберг, Кано, Панаццоло), индекс Пуанкаре--Хопфа, центральные многообразия, нормальная гиперболичность (Фенихель), метод раздутий Дюмортье--Руссари\\
Небесная механика & раздутие Мак-Ги, коллизионное многообразие, эффективные потенциалы\\
Топология & линки и коническая структура, эйлерова характеристика, принцип Важевского, индекс Конли\\
Вариационные методы & принцип Мопертюи--Якоби, Хопф--Ринов, тормозные орбиты (Болотин), Хагедорн--Мавен\\
Вещественная алгебраическая геометрия форм & неотрицательные тернарные квартики (Гильберт), структура нулей\\
Теория возмущений & адиабатические инварианты (интуиция о <<зеркальном>> эффекте и о том, почему гладкие контрпримеры возможны)\\
\bottomrule
\end{tabularx}
\end{center}

%=====================================================================
\section*{Литература}
\addcontentsline{toc}{section}{Литература}
\small
\begin{itemize}[leftmargin=*,itemsep=1pt]
\item[{[Ar]}] В.\,И.~Арнольд (ред.), \emph{Задачи Арнольда}, задача 1971-4 и комментарии.
\item[{[AR17]}] G.~Allaire, J.~Rauch, Instability of dielectrics and conductors in electrostatic fields, ARMA 224 (2017).
\item[{[BB00]}] M.\,L.~Bertotti, S.\,V.~Bolotin, On the influence of the kinetic energy on the stability of equilibria of natural Lagrangian systems, ARMA 152 (2000).
\item[{[BMP21]}] J.\,M.~Burgos, E.~Maderna, M.~Paternain, On the Lyapunov instability in Newtonian dynamics, Nonlinearity 34 (2021); \href{https://arxiv.org/abs/1906.11962}{arXiv:1906.11962}.
\item[{[BP22]}] J.\,M.~Burgos, M.~Paternain, On the Lagrange--Dirichlet converse in dimension three, \href{https://arxiv.org/abs/2208.01139}{arXiv:2208.01139}.
\item[{[Br98]}] M.~Brunella, Instability of equilibria in dimension three, Ann. Inst. Fourier 48:5 (1998), 1345--1357; \href{http://www.numdam.org/item/AIF_1998__48_5_1345_0/}{numdam}.
\item[{[Bu21]}] J.\,M.~Burgos, A proof of the Palamodov's total instability conjecture, \href{https://arxiv.org/abs/2108.05829}{arXiv:2108.05829} (отозвана).
\item[{[Bu25]}] J.\,M.~Burgos, McGehee blowup for Lagrangian systems and instability of equilibria, \href{https://arxiv.org/abs/2506.19135}{arXiv:2506.19135}; J. Differential Equations 454 (2026).
\item[{[Ch52]}] Н.\,Г.~Четаев, О неустойчивости равновесия в некоторых случаях, когда функция сил не есть максимум, ПММ 16 (1952).
\item[{[Ha71]}] P.~Hagedorn, Die Umkehrung der Stabilitätssätze von Lagrange--Dirichlet und Routh, ARMA 42 (1971).
\item[{[KP82]}] В.\,В.~Козлов, В.\,П.~Паламодов, Об асимптотических решениях уравнений классической механики, ДАН СССР 263:2 (1982).
\item[{[Ko81]}] В.\,В.~Козлов, О неустойчивости равновесия в потенциальном поле, УМН 36:3 (1981) (англ.: Russian Math. Surveys 36 (1981), 238--239).
\item[{[Ko82]}] В.\,В.~Козлов, Асимптотические решения уравнений классической механики, ПММ 46:4 (1982); Гипотеза о существовании асимптотических движений в классической механике, Функц. анализ 16:4 (1982).
\item[{[Ko86]}] В.\,В.~Козлов, Асимптотические движения и проблема обращения теоремы Лагранжа--Дирихле, ПММ 50:6 (1986).
\item[{[KT18]}] V.\,V.~Kozlov, D.\,V.~Treschev, Instability, asymptotic trajectories and dimension of the phase space, Mosc. Math. J. 18:4 (2018).
\item[{[Ku89]}] А.\,Н.~Кузнецов, О существовании входящих в особую точку решений автономной системы, обладающей формальным решением, Функц. анализ 23:4 (1989).
\item[{[LP82]}] M.~Laloy, K.~Peiffer, On the instability of equilibrium when the potential has a non-strict local minimum, ARMA 78 (1982).
\item[{[MN89]}] V.~Moauro, P.~Negrini, On the inversion of the Lagrange--Dirichlet theorem, Diff. Int. Eq. 2 (1989).
\item[{[Pa77]}] В.\,П.~Паламодов, Об устойчивости равновесия в потенциальном поле, Функц. анализ и его прил. 11:4 (1977), 42--55.
\item[{[Pa95]}] V.\,P.~Palamodov, Stability of motion and algebraic geometry, AMS Transl. (2) 168 (1995), 5--20.
\item[{[Pa20]}] В.\,П.~Паламодов, Об обращении теоремы Лагранжа--Дирихле и неустойчивости консервативных систем, УМН 75:3 (2020), 107--122; \href{https://www.mathnet.ru/eng/rm9945}{mathnet}. \textbf{Лучший обзор по теме на русском.}
\item[{[Ta80a]}] S.\,D.~Taliaferro, An inversion of the Lagrange--Dirichlet stability theorem, ARMA 73 (1980).
\item[{[Ta80b]}] S.\,D.~Taliaferro, Stability for two dimensional analytic potentials, JDE 35 (1980), 248--265.
\item[{[Ta90]}] S.\,D.~Taliaferro, Instability of an equilibrium in a potential field, ARMA 109 (1990).
\item[{[Ur20]}] A.\,J.~Ureña, To what extent are unstable the maxima of the potential?, Ann. Mat. Pura Appl. 199 (2020).
\item[{[Wi93]}] Z.~Wiener, Instability with two zero frequencies, JDE 103 (1993).
\item[] Также: Кано (1987), Панаццоло (Acta Math. 2006), Маккуиллан--Панаццоло (2013) --- редукция особенностей полей в размерности~3; Kurdyka--Mostowski--Parusiński, Ann. of Math. 152 (2000) --- гипотеза Тома о градиенте; K.~Saito, Invent. Math. 14 (1971); Bochnak--Łojasiewicz (1971); Гарсия--Таль, JDE 194 (2003); Козлов--Фурта, \emph{Асимптотики решений сильно нелинейных систем}.
\end{itemize}

\end{document}
